\documentclass{article}

\usepackage[T1]{fontenc}

\usepackage{booktabs,graphicx,siunitx}

\begin{document}

\section{Wilcoxon signed-rank}

\begin{tabular}{llll}
\toprule
Sign & N & Mean rank & Sum of ranks \\
\midrule
Positive & 40 & 20.50 & 820.00 \\
Negative & 0 & --- & --- \\
Ties & 0 & --- & --- \\
\bottomrule
\end{tabular}

\begin{tabular}{lllll}
\toprule
V / W & Z & p & r [95\% CI] & Median diff [95\% CI] \\
\midrule
820.00 & 5.51 & $<$.001 & 1.00 [1.00, 1.00] & 72.24 [68.63, 75.10] \\
\bottomrule
\end{tabular}

\includegraphics{figures/01_wilcoxon-signed-rank/figure_difference.png}

The nonparametric counterpart to the paired t-test. It ranks the size of each paired change; a p below alpha means a systematic shift between the two conditions, with r as the effect size. The shift maps to the median difference only when the within-pair differences are roughly symmetric; the sign of the median difference (or Hodges--Lehmann estimate) gives the direction. Your significance threshold ($\alpha$) is 0.05. Method: Wilcoxon signed rank exact test.

A Wilcoxon signed-rank test gave V=820.00, Z=5.51, p $<$ .001, r=1.00 [1.00, 1.00].

Statistical basis: Wilcoxon signed-rank test (Wilcoxon, F. (1945). "Individual comparisons by ranking methods." Biometrics Bulletin, 1(6), 80-83.); Matched-pairs rank-biserial effect size (Ben-Shachar MS, Lüdecke D, Makowski D (2020). "effectsize: Estimation of Effect Size Indices and Standardized Parameters." Journal of Open Source Software, 5(56), 2815.).

\end{document}